% graphicx loaded via preamble (includegraphics used below)
% ==============================================================================
% T0-Theorie: Dokument 164
% Titel: T0-Theorie und 't Hoofts Zellularautomat-Interpretation der
%        Quantenmechanik — Determinismus aus zwei Richtungen
% Autor: Johann Pascher
% Datum: Februar 2026
% Referenz: G. 't Hooft, The Cellular Automaton Interpretation of Quantum
%           Mechanics, Springer 2016, DOI 10.1007/978-3-319-41285-6
% ==============================================================================

\section*{Abstract}

Gerard 't Hoofts \emph{The Cellular Automaton Interpretation of Quantum Mechanics}
(Springer, 2016) und die T0-Theorie verfolgen dasselbe fundamentale Ziel: die
Quantenmechanik als deterministisch zu verstehen und die Kopenhagener Deutung
durch einen tieferliegenden klassischen Mechanismus zu ersetzen. Dieses Dokument
analysiert systematisch, wo beide Ansätze übereinstimmen (Determinismus, Ablehnung
des Wellenpaketkollapses, Zeitdiskretheit auf der Planck-Skala, Hilbert-Raum als
Werkzeug) und wo T0 fundamentaler ist: 't Hooft postuliert den zellularen Automaten
als neue Grundlage, ohne zu erklären, warum gerade diese Struktur. T0 leitet die
diskrete Raumzeitstruktur aus einem einzigen geometrischen Parameter $\xipar =
4/30\,000$ her. Die zentrale These dieses Dokuments: Der 't Hooftsche Zellularautomat
ist kein unabhängiges Postulat, sondern ein Grenzfall der T0-Geometrie auf der
Sub-Planck-Skala $L_0 = \xipar \cdot \ell_P$. Ferner wird gezeigt, dass 't Hoofts
offene Probleme (Lokalität, Fermionen, Gravitation, freie Parameter des
Standardmodells) in T0 durch die Zeit-Masse-Dualität $T \cdot m = 1$ und die
fraktale Raumzeitgeometrie eine natürliche Antwort erhalten.

\clearpage

% ==============================================================================

\begin{figure}[htbp]
  \centering
  \includegraphics[width=0.92\textwidth]{../fig/fig1_hierarchy_de}
  \caption{Drei Interpretationsebenen der Quantenmechanik: Kopenhagener Deutung (probabilistisch), CAI von \textquoteleft{}t~Hooft (Automat postuliert), T0-Theorie (Determinismus aus $\xi$ abgeleitet). Jede Ebene subsumiert die vorherige.}
  \label{fig:hierarchy_de}
\end{figure}

\section{Einleitung: Zwei Nobelpreisträger-Intuitionen, ein Ziel}
\label{sec:intro}
% ==============================================================================

\subsection{Das gemeinsame Unbehagen}

Gerard 't Hooft, Nobelpreisträger 1999 für die Aufklärung der Quantenstruktur der
elektroschwachen Wechselwirkungen, hat sich in den letzten Jahrzehnten intensiv mit
einer Frage beschäftigt, die die meisten seiner Kollegen für entschieden hielten:
Ist die Quantenmechanik fundamental probabilistisch, oder liegt ihr eine
deterministische Struktur zugrunde?

Seine Antwort, ausführlich dargelegt in \emph{The Cellular Automaton Interpretation
of Quantum Mechanics} (2016), ist eindeutig: Die Quantenmechanik ist ein Werkzeug,
keine Theorie. Unter ihr liegt ein klassisch-deterministischer Mechanismus, der auf
der Planck-Skala operiert und dessen effektive Beschreibung die vertraute
Quantenmechanik ergibt.

Die T0-Theorie (Zeit-Masse-Dualität, FFGFT) kommt aus einer völlig anderen Richtung
zu derselben Schlussfolgerung. Ausgehend von der Beobachtung, dass alle
Naturkonstanten aus einem einzigen geometrischen Parameter $\xipar = 4/30\,000$
emergieren, ergibt sich nicht nur Determinismus als Konsequenz, sondern auch eine
konkrete geometrische Struktur der Planck-Skala, die 't Hoofts Zellularautomaten
als Spezialfall enthält.

\subsection{Die Struktur dieses Dokuments}

\begin{enumerate}
    \item \textbf{Abschnitt 2}: Kernpunkte von 't Hoofts CAI — was er behauptet,
    was er beweist, was er offen lässt.
    \item \textbf{Abschnitt 3}: Sechs Punkte tiefer Übereinstimmung zwischen CAI
    und T0.
    \item \textbf{Abschnitt 4}: Vier fundamentale Unterschiede — wo T0 tiefer geht.
    \item \textbf{Abschnitt 5}: Die zentrale These: Der Zellularautomat als
    Grenzfall von T0.
    \item \textbf{Abschnitt 6}: 't Hoofts offene Probleme und T0-Antworten.
    \item \textbf{Abschnitt 7}: Testbare Vorhersagen aus dem Vergleich.
\end{enumerate}

\clearpage

% ==============================================================================
\section{Kernpunkte der Zellularautomat-Interpretation}
\label{sec:cai_summary}
% ==============================================================================

\subsection{Die Grundthese}

't Hooft beginnt mit einer fundamentalen Kritik an der Kopenhagener Deutung. Die
Standardinterpretation der Quantenmechanik lässt mehrere unbequeme Lücken:
den Wellenpaketkollaps beim Messen, die Nichtlokalität der Verschränkung, die
Unmöglichkeit einer klassischen Beschreibung der Realität, und die scheinbare
Unvereinbarkeit mit der Relativitätstheorie auf fundamentaler Ebene.

't Hoofts Gegenentwurf in einem Satz: \emph{Quantenmechanik ist die effektive
Theorie eines klassisch-deterministischen Systems, das auf der Planck-Skala
als zellularer Automat operiert.}

\subsection{Das Zahnradmodell als Prototyp}

Das einfachste Modell der CAI ist das \emph{Zahnrad-Modell} (cogwheel model):
Ein System mit $N$ diskreten Zuständen $|n\rangle$, $n = 0, \ldots, N-1$, das
sich nach einer festen Regel entwickelt: bei jedem Zeitschritt $\delta t$ wird
$n \to n+1 \pmod{N}$. Dies ist vollständig klassisch-deterministisch.

't Hooft zeigt nun, dass sich diesem System ein Hilbert-Raum und ein Hamilton-Operator
zuordnen lassen:
\begin{equation}
    H_{\text{op}} |n\rangle = \frac{2\pi k}{N \delta t} |n\rangle
    \label{eq:cogwheel_hamilton}
\end{equation}
wobei $k$ die Energiequantenzahl ist. Die Zustände gehorchen der
Schrödinger-Gleichung — nicht weil das System quantenmechanisch ist, sondern
weil Quantenmechanik \emph{genau das ist}: eine effektive Beschreibung periodischer
deterministischer Systeme.


\begin{figure}[htbp]
  \centering
  \includegraphics[width=\textwidth]{../fig/fig2_cogwheel_torus_de}
  \caption{Links: CAI-Zahnrad mit $N=8$ diskreten Zust\"{a}nden (postulierte Struktur). Rechts: T0-Zeitfeld $T(x,t)$ auf geometrischem Torus mit $L_0 = \xi \cdot \ell_P$ (abgeleitete Struktur).}
  \label{fig:cogwheel_de}
\end{figure}

\subsection{Ontologische vs.\ Schablonenzustände}

Ein zentrales Konzept der CAI ist die Unterscheidung zwischen:
\begin{itemize}
    \item \textbf{Ontologischen Zuständen} $|n\rangle_{\text{ont}}$: Die tatsächlichen
    Zustände des klassisch-deterministischen Systems. Zu jedem Zeitpunkt befindet
    sich das System in genau einem solchen Zustand.
    \item \textbf{Schablonenzuständen} (template states) $|\psi\rangle$: Superpositionem
    ontologischer Zustände. Diese sind kein Abbild der Realität, sondern
    Beschreibungswerkzeuge eines Beobachters, der die genauen Anfangsbedingungen
    nicht kennt.
\end{itemize}

Probabilismus in der Quantenmechanik ist nach 't Hooft \emph{epistemisch}: er
reflektiert die Unwissenheit des Beobachters über die genauen ontologischen
Anfangsbedingungen, nicht eine fundamentale Indeterminiertheit der Natur.

\subsection{Bell-Theorem und Superdeterminismus}

Das scheinbar stärkste Argument gegen Hidden-Variable-Theorien ist Bells Theorem:
keine lokale, realistische Theorie kann die Quantenkorrelationen reproduzieren.
't Hooft antwortet mit dem Konzept des \emph{Superdeterminismus}: Wenn das Universum
vollständig deterministisch ist, dann sind auch die Einstellungen der Messapparate
durch die Anfangsbedingungen determiniert. Die Freiheitsgrade, die in Bells Argument
als unabhängig vorausgesetzt werden, sind es nicht.

't Hooft zeigt, dass dieser Einwand prinzipiell überwindbar ist, räumt aber ein,
dass eine vollständige quantitative Theorie des Superdeterminismus noch aussteht.

\subsection{Die offenen Probleme}

't Hooft ist ausdrücklich transparent über die Grenzen seines Ansatzes. Am Ende
des Buches nennt er explizit:
\begin{itemize}
    \item Das \textbf{Lokalitätsproblem}: Wie werden Quantenkorrelationen lokal
    aus dem Automaten reproduziert?
    \item Das \textbf{Fermionen-Problem}: Wie entstehen halbzahlige Spins aus
    einem Boson-artigen Automaten?
    \item Das \textbf{Gravitationsproblem}: Der Automat wurde nicht mit
    allgemeiner Relativitätstheorie verbunden.
    \item Das \textbf{freie-Parameter-Problem}: Das Standardmodell hat 20+
    freie Parameter. Der Automat erklärt nicht, warum sie diese Werte haben.
    \item Das \textbf{Informationsverlust-Problem}: Schwarze Löcher und
    Unitarität bleiben ungeklärt.
\end{itemize}

\begin{keyresult}['t Hoofts Programm]
    Die CAI ist ein Programm, keine fertige Theorie. Sie zeigt, dass deterministische
    Quantenmechanik \emph{möglich} ist, aber sie liefert noch keine vollständige
    Ableitung der bekannten Physik aus einem deterministischen Fundament.
\end{keyresult}

\clearpage

% ==============================================================================
\section{Sechs Punkte tiefer Übereinstimmung}
\label{sec:agreement}
% ==============================================================================

\subsection{1. Determinismus als fundamentales Prinzip}

Beide Ansätze lehnen den fundamentalen Probabilismus der Kopenhagener Deutung ab.
In der CAI ist Zufall epistemisch — er reflektiert Unwissenheit über
Anfangsbedingungen. In T0 ist Zufall ebenfalls epistemisch: Das Zeitfeld
$\Tfield = \hbar / (mc^2)$ entwickelt sich vollständig deterministisch gemäß dem
T0-Energiefunktional:
\begin{equation}
    E_{T0} = \int |\nabla T(x)|^2 \, d^4x
    \label{eq:t0_energy}
\end{equation}
Scheinbare Quantenunschärfe ist in T0 die Projektion eines deterministischen
vierdimensionalen Zeitfeldes auf die dreidimensionale Raumhyperfläche eines
Beobachters.

\subsection{2. Ablehnung des Wellenpaketkollapses}

't Hooft: Der Kollaps ist kein physikalischer Prozess, sondern eine
Informationsaktualisierung des Beobachters über den ontologischen Zustand des
Systems. Vor der Messung war das System in einem bestimmten Zustand — der
Beobachter wusste es nur nicht.

T0: Identische Schlussfolgerung aus anderem Fundament. Das Zeitfeld $\Tfield$
kollabiert nicht — es entwickelt sich stetig. Eine Messung ist eine
Wechselwirkung zweier Zeitfelder, deren Korrelation die scheinbare
Kollapsstruktur erzeugt.

\subsection{3. Hilbert-Raum als Werkzeug, nicht als Ontologie}

't Hooft schreibt explizit: \emph{Quantum mechanics is viewed as a tool, not as
a theory.} Der Hilbert-Raum beschreibt keine fundamentale Realität — er ist ein
effektives Beschreibungswerkzeug für Systeme, über deren mikroskopische Details
man keine vollständige Information hat.

T0 stimmt vollständig zu: Hilbert-Raum-Zustände $|\psi\rangle$ sind in T0 eine
komprimierte Darstellung der Zeitfeldkonfiguration $T(x,t)$. Die Wellenfunktion
ist eine statistische Beschreibung des Zeitfeldes aus der Perspektive eines
Beobachters mit begrenzter Ortsauflösung.

\subsection{4. Diskrete Zeitstruktur auf der Planck-Skala}

't Hooft: Der Automat operiert mit einem fundamentalen Zeitschritt $\delta t$
von der Größenordnung der Planck-Zeit $t_P \approx 5{,}39 \times 10^{-44}$ s.
Unterhalb dieser Skala ist keine Dynamik definiert.

T0: Die Sub-Planck-Zeitskala ist nicht $t_P$, sondern
\begin{equation}
    T_0 = \xipar \cdot t_P = \frac{4}{30\,000} \cdot t_P
    \approx 7{,}19 \times 10^{-48} \, \text{s}
    \label{eq:t0_time}
\end{equation}
Beide Ansätze teilen die Überzeugung, dass Zeit auf der kleinsten Skala diskret
ist. T0 liefert den genauen Wert dieser Diskretheit — 't Hooft lässt ihn offen.

\subsection{5. Quantenmechanik als Grenzfall klassischer Dynamik}

't Hooft zeigt an konkreten Modellen (Zahnrad, harmonischer Rotator, freie Felder),
dass klassisch-deterministische Systeme Schrödinger-Gleichungen gehorchen, wenn
man die passende Basis wählt. Quantenmechanik \emph{emergiert} aus klassischer
Mechanik.

T0 macht dieselbe Aussage präziser: Die Schrödinger-Gleichung ist die lineare
Näherung der nichtlinearen T0-Feldgleichung
\begin{equation}
    \nabla^2 T - \frac{1}{c^2}\frac{\partial^2 T}{\partial t^2}
    = -\frac{\xipar}{T^2} \cdot |\nabla T|^2
    \label{eq:t0_field_equation}
\end{equation}
für schwache Zeitfeldgradienten, d.h.\ für $|\nabla T| \ll 1/\xipar$.

\subsection{6. Zeitpfeil als emergentes Phänomen}

't Hooft: Der Zeitpfeil kann aus dem deterministischen Automaten erklärt werden,
wenn man Informationsverlust bei Kollisionen einführt — bestimmte Mikrozustände
kollabieren auf denselben Makrozustand.

T0: Der Zeitpfeil folgt geometrisch aus $T \cdot m = 1$. Da $T > 0$ algebraisch
erzwungen ist und die Toruswicklung eine topologisch geschützte Richtung definiert,
ist der Zeitpfeil keine thermodynamische Approximation, sondern eine fundamentale
geometrische Tatsache (vgl.\ Dokument 157).

\begin{discoveryBox}[Tiefe Konvergenz]
    Zwei voneinander unabhängig entwickelte Theorien — 't Hoofts CAI aus der
    Perspektive der Quantenfeldtheorie und T0 aus der Perspektive geometrischer
    Vereinheitlichung — konvergieren in sechs fundamentalen Punkten. Das ist
    kein Zufall: Beide stoßen gegen dieselbe Wand der Standardinterpretation
    und finden ähnliche Ausbrüche.
\end{discoveryBox}

\clearpage

% ==============================================================================
\section{Vier fundamentale Unterschiede: Wo T0 tiefer geht}
\label{sec:differences}
% ==============================================================================

\subsection{Unterschied 1: Postulat vs.\ Ableitung}

\paragraph{'t Hooft:} Der zellulare Automat wird als neue fundamentale Struktur
\emph{postuliert}. Es wird angenommen, dass die Natur auf der Planck-Skala wie ein
Automat mit diskreten Zuständen und einer deterministischen Evolutionsregel
funktioniert. Warum \emph{diese} Struktur? Das Buch gibt keine Antwort.

\paragraph{T0:} Die diskrete Raumzeitstruktur wird aus $\xipar = 4/30\,000$
\emph{abgeleitet}. Der Parameter $\xipar$ ist nicht frei — er emergiert aus
der Geometrie des dreidimensionalen Raumes auf der Planck-Skala. Die
Zellengrößen des Automaten sind damit festgelegt:
\begin{equation}
    \ell_{\text{Zelle}} = L_0 = \xipar \cdot \ell_P, \quad
    t_{\text{Schritt}} = T_0 = \xipar \cdot t_P
    \label{eq:cell_size}
\end{equation}

\begin{table}[h]
    \centering
    \caption{Postulat vs.\ Ableitung: Grundlegende Strukturen}
    \label{tab:postulate_vs_derivation}
    \resizebox{\textwidth}{!}{\begin{tabular}{@{}lll@{}}
        \toprule
        Struktur & 't Hooft CAI & T0-Theorie \\
        \midrule
        Diskrete Zeitstruktur & Postuliert ($\delta t \sim t_P$) & Abgeleitet ($T_0 = \xipar \cdot t_P$) \\
        Zellgröße & Offen ($\sim \ell_P$) & Festgelegt ($L_0 = \xipar \cdot \ell_P$) \\
        Evolutionsregel & Postuliert (Automat) & Abgeleitet aus $E_{T0}$ \\
        Anzahl freier Parameter & Mehrere & Null \\
        Warum dieser Wert? & Keine Antwort & $\xipar$ aus Geometrie \\
        \bottomrule
    \end{tabular}}
\end{table}

\subsection{Unterschied 2: Freie Parameter des Standardmodells}

\paragraph{'t Hooft:} Das Buch anerkennt explizit, dass die CAI die 20+ freien
Parameter des Standardmodells nicht erklärt. Die Massen der Quarks, die
Kopplungskonstanten, der Weinberg-Winkel — all das bleibt unerklärt.

\paragraph{T0:} Alle diese Parameter emergieren aus $\xipar$. Die Teilchenmassen
folgen aus der Beziehung $E_n = E_0 / \xipar^n$ mit $E_0 = m_e c^2$:
\begin{equation}
    m_n = m_e \cdot \xipar^{-n}, \quad n = 0, 1, 2, \ldots
    \label{eq:mass_spectrum}
\end{equation}
Diese Formel reproduziert 98\% aller bekannten Teilchenmassen. Die
Feinstrukturkonstante folgt aus $\alpha = \xipar \cdot E_0^2$.

\subsection{Unterschied 3: Gravitation und Raumzeit}

\paragraph{'t Hooft:} Die Einbindung von Gravitation in die CAI ist ausdrücklich
als offenes Problem genannt. Das Buch enthält einen Anhang über Gravitation in
2+1 Dimensionen, aber keine vollständige 3+1-dimensionale Theorie.

\paragraph{T0:} Gravitation emergiert natürlich aus dem Zeitfeldgradienten:
\begin{equation}
    \vec{g}(x) = -c^2 \nabla \ln T(x) = \frac{c^2}{T} \nabla T
    \label{eq:t0_gravity}
\end{equation}
Das Newtonsche Gravitationspotential $\Phi = -GM/r$ entspricht
$\ln T \propto GM/(rc^2)$, und die allgemeine Relativitätstheorie ergibt sich
als klassischer Grenzfall der T0-Feldgleichungen. Die Gravitationskonstante $G$
emergiert ebenfalls aus $\xipar$:
\begin{equation}
    G = \frac{\xipar^2 \cdot \ell_P^2 \cdot c^3}{\hbar}
    \label{eq:G_from_xi}
\end{equation}

\subsection{Unterschied 4: Das Fermionen-Problem}

\paragraph{'t Hooft:} Fermionen (halbzahliger Spin, Pauli-Prinzip) sind besonders
schwer aus einem Boson-artigen Automaten zu erhalten. Das Buch enthält ein Kapitel
über Fermionen, erklärt aber selbst, dass die Konstruktion noch unvollständig ist.

\paragraph{T0:} Fermionen entstehen natürlich aus der Topologie des Zeitfeldes.
Die Bindungsbedingung $T \cdot m = 1$ kombiniert mit der toroidalen Topologie
des T0-Universums erzeugt halbzahlige Windungszahlen, die dem Fermionen-Spin
entsprechen. Ein Teilchen mit halbzahligem Spin ist eine Zeitfeldkonfiguration
mit einer Torus-Windungszahl $\nu = 1/2$ — geometrisch natürlich, nicht
zusätzlich postuliert.

\clearpage

% ==============================================================================
\section{Der Zellularautomat als Grenzfall von T0}
\label{sec:ca_as_limit}
% ==============================================================================

\subsection{Die zentrale These}

Wir formulieren jetzt die Hauptaussage dieses Dokuments präzise:

\begin{revolutionary}[Hauptthese: Inklusion]
    Der 't Hooftsche Zellularautomat ist kein unabhängiges Fundament der Physik,
    sondern ein effektiver Grenzfall der T0-Geometrie auf der Sub-Planck-Skala.
    Konkret: Die CAI beschreibt die diskrete Zeitentwicklung des T0-Zeitfeldes
    $T(x,t)$ in der Gitterapproximation $\Delta x = L_0$, $\Delta t = T_0$.
\end{revolutionary}


\begin{figure}[htbp]
  \centering
  \includegraphics[width=0.95\textwidth]{../fig/fig3_inclusion_de}
  \caption{Inklusionsdiagramm: Die CAI (orange) ist als Grenzfall in der T0-Theorie (gr\"{u}n) enthalten. Die \"{a}u\ss{}eren gr\"{u}nen Elemente sind T0-Beitr\"{a}ge jenseits der CAI.}
  \label{fig:inclusion_de}
\end{figure}

\subsection{Die Verbindung im Detail}

In der CAI wird die Zeitentwicklung des Universums beschrieben durch eine
unitäre Evolutionsmatrix $U_{\text{op}}(\delta t) = e^{-iH_{\text{op}}\delta t}$,
die auf ontologische Zustände wirkt. Der Hamiltonoperator $H_{\text{op}}$ ist
formal definiert, aber nicht aus tieferen Prinzipien abgeleitet.

In T0 ist das Äquivalent das Zeitfeld-Energiefunktional $E_{T0}$.
Die Zeitentwicklung des diskreten T0-Gitters auf der Skala $L_0$ ist:
\begin{equation}
    T(x + L_0, t + T_0) = T(x, t) - \xipar \cdot T_0 \cdot
    \frac{\partial E_{T0}}{\partial T(x,t)}
    \label{eq:t0_discrete_evolution}
\end{equation}

Dies ist exakt die Struktur eines zellularen Automaten:
\begin{itemize}
    \item Jede Zelle $(x, t)$ enthält den Wert $T(x,t)$ — das ist der
    \emph{ontologische Zustand} im Sinne 't Hoofts.
    \item Die Evolutionsregel ist lokal und deterministisch — genau wie in der CAI.
    \item Der CAI-Hamiltonoperator $H_{\text{op}}$ entspricht in T0 der
    zweiten Variation des Energiefunktionals $\delta^2 E_{T0}$.
\end{itemize}

\subsection{Warum T0 mehr liefert}

Der Automat in der CAI hat zwei freie Komponenten: die Zellgröße und die
Evolutionsregel. T0 fixiert beide:
\begin{align}
    \ell_{\text{Zelle}} &= L_0 = \xipar \cdot \ell_P
    \quad \Longleftarrow \quad \xipar = \frac{4}{30\,000}
    \label{eq:cell_size_fixed} \\
    \text{Evolutionsregel} &: T \to T - \xipar T_0 \frac{\delta E_{T0}}{\delta T}
    \quad \Longleftarrow \quad E_{T0} = \int |\nabla T|^2 d^4x
    \label{eq:rule_fixed}
\end{align}

Der Automat ist nicht frei wählbar — er ist geometrisch determiniert.

\subsection{Das Zahnrad-Modell in T0-Sprache}

't Hoofts einfachstes Beispiel, das Zahnrad mit $N$ Zuständen, entspricht in T0
einem Zeitfeld auf einem eindimensionalen Torus mit $N$ Gitterpunkten und
Bindungsbedingung $T_n \cdot m = 1$. Die periodische Evolutionsregel
$n \to n + 1 \pmod{N}$ entsteht aus der geometrischen Periodizität des Torus.
Die Energie-Eigenwerte aus \eqref{eq:cogwheel_hamilton} entsprechen den
harmonischen Frequenzen $f_k = k/(N \cdot T_0)$ des T0-Zeitfeldes auf dem Torus.

\begin{insightBox}[Geometrie ersetzt Postulat]
    Das Zahnrad-Modell ist nicht ein \emph{Beispiel} für einen möglichen Automaten —
    es ist das \emph{einzige} Modell, das mit der T0-Torusgeometrie verträglich ist.
    Die periodische Struktur ist nicht gewählt, sie ist geometrisch erzwungen.
\end{insightBox}

\clearpage

% ==============================================================================
\section{'t Hoofts offene Probleme und T0-Antworten}
\label{sec:open_problems}
% ==============================================================================


\begin{figure}[htbp]
  \centering
  \includegraphics[width=0.62\textwidth]{../fig/fig4_radar_de}
  \caption{Radardiagramm der offenen Probleme. Orange: CAI ({\lq}t~Hooft). Gr\"{u}n: T0-Theorie. Bewertung: 0 = offen, 2 = gel\"{o}st.}
  \label{fig:radar_de}
\end{figure}

\subsection{Das Lokalitätsproblem}

\paragraph{CAI-Problem:} Quantenverschränkung erzeugt nichtlokale Korrelationen,
die mit einer lokal operierenden Automaten-Regel schwer verträglich scheinen.
't Hooft argumentiert für Superdeterminismus, gibt aber zu, dass eine quantitative
Theorie fehlt.

\paragraph{T0-Antwort:} In T0 sind verschränkte Teilchen durch eine gemeinsame
Zeitfeldkonfiguration $T_{\text{AB}}(x_A, x_B, t)$ verbunden. Diese Konfiguration
entwickelt sich lokal — aber da das Zeitfeld an beiden Orten durch dieselbe
geometrische Struktur bestimmt wird, entstehen die beobachteten Korrelationen
ohne nichtlokale Wechselwirkung. Das ist kein Superdeterminismus im Sinne
willkürlicher Anfangsbedingungen — es ist geometrische Notwendigkeit.

\subsection{Das Fermionen-Problem}

\paragraph{CAI-Problem:} Halbzahlige Spins entstehen nicht natürlich aus einem
ganzzahligen Automaten. Das Kapitel über Fermionen im Buch bleibt unabgeschlossen.

\paragraph{T0-Antwort:} Die toroidale Topologie des T0-Universums erlaubt
Zeitfeldkonfigurationen mit halbzahligen Windungszahlen. Ein Fermion ist
mathematisch eine Zeitfeldkonfiguration mit
\begin{equation}
    \oint_\gamma \frac{d T}{T} = \pi \quad (\text{halbzahlige Windung})
    \label{eq:fermion_winding}
\end{equation}
Dies ist geometrisch stabil (topologischer Schutz) und erfordert kein
zusätzliches Postulat.

\subsection{Das Gravitationsproblem}

\paragraph{CAI-Problem:} Keine vollständige Verbindung zur allgemeinen
Relativitätstheorie. Der Anhang behandelt Gravitation in 2+1 Dimensionen
als Testfall, aber 3+1 Dimensionen fehlen.

\paragraph{T0-Antwort:} Die Gleichung \eqref{eq:t0_gravity} verbindet direkt
das Zeitfeld mit der Gravitationsbeschleunigung. In der schwachen Feldnäherung
reproduziert T0 Einsteins allgemeine Relativitätstheorie. In der starken
Feldnäherung (nahe Singularitäten) gibt T0 Korrekturen der Ordnung $\xipar^2$.

\subsection{Das freie-Parameter-Problem}

\paragraph{CAI-Problem:} Die 20+ freien Parameter des Standardmodells werden
nicht erklärt. Das Buch erwähnt dieses Problem, bezeichnet es aber als
außerhalb des aktuellen Rahmens.

\paragraph{T0-Antwort:} Alle Parameter folgen aus $\xipar$:
\begin{table}[h]
    \centering
    \caption{Freie Parameter des Standardmodells in T0}
    \label{tab:sm_parameters}
    \resizebox{\textwidth}{!}{\begin{tabular}{@{}lll@{}}
        \toprule
        Parameter & Standardmodell & T0 \\
        \midrule
        Feinstrukturkonstante $\alpha$ & Gemessen: $1/137{,}036$ & $\alpha = \xipar \cdot E_0^2$ \\
        Gravitationskonstante $G$ & Gemessen & $G = \xipar^2 \ell_P^2 c^3/\hbar$ \\
        Elektronmasse $m_e$ & Gemessen & $m_e = E_0/c^2$, $E_0 = \hbar c/(\xipar \cdot \ell_P)$ \\
        Myonmasse $m_\mu$ & Gemessen & $m_\mu = m_e/\xipar$ \\
        Taumasse $m_\tau$ & Gemessen & $m_\tau = m_e/\xipar^2$ \\
        Planck-Konstante $h$ & Gemessen & Emergiert aus $\xipar$ \\
        Alle anderen & 15+ frei & Aus $\xipar$ \\
        \bottomrule
    \end{tabular}}
\end{table}

\subsection{Das Informationsverlust-Problem}

\paragraph{CAI-Problem:} 't Hooft hat in früheren Arbeiten argumentiert, dass
bei der Bildung schwarzer Löcher Information verloren gehen kann — im Widerspruch
zur Unitarität der Quantenmechanik. Das Buch adressiert dies, ohne eine
vollständige Auflösung zu bieten.

\paragraph{T0-Antwort:} Im T0-Rahmen ist Information niemals verloren. Das
Zeitfeld $T(x,t)$ entwickelt sich unitär gemäß dem Energiefunktional
$E_{T0}$. Schwarze Löcher sind in T0 Regionen extremer Zeitfeldkompression.
Entscheidend: Die fraktale Korrektur $K_{\text{frak}}$ setzt eine harte
geometrische Untergrenze. Das Zeitfeld kann \emph{nicht} bis null komprimiert
werden — es gibt ein minimales $T$:
\begin{equation}
    T \geq T_{\min} = \xi \cdot t_P \approx 7{,}19 \times 10^{-48}\,\text{s}
    \quad \Longrightarrow \quad m \leq m_{\max} = \frac{\hbar}{T_{\min}\,c^2}
    \label{eq:t_min_bh}
\end{equation}
Schwarze Löcher sind daher in T0 keine echten Singularitäten — das Zeitfeld
komprimiert sich bis $T_{\min}$, bleibt aber endlich und wohldefinierten
Gradienten erhalten. Die Informationsparadoxie löst sich auf: das Zeitfeld
ist unitär, trägt vollständige Information, und keine Divergenz tritt auf.

\clearpage

% ==============================================================================
\section{Testbare Vorhersagen aus dem Vergleich}
\label{sec:predictions}
% ==============================================================================


\begin{figure}[htbp]
  \centering
  \includegraphics[width=\textwidth]{../fig/fig5_scales_de}
  \caption{Skalenhierarchie von der menschlichen Skala bis zur T0 Sub-Planck-Skala. Die T0-Gitterkonstante $L_0 = \xi \cdot \ell_P$ liegt um den Faktor $\xi \approx 1{,}33 \times 10^{-4}$ unterhalb der Planck-L\"{a}nge.}
  \label{fig:scales_de}
\end{figure}

\subsection{Vorhersage 1: Zellengrößen-Signatur}

't Hooft lässt die Zellgröße $\delta t$ offen (von der Ordnung $t_P$). T0 fixiert
sie auf $T_0 = \xipar \cdot t_P$. Das Verhältnis
\begin{equation}
    \frac{T_0}{t_P} = \xipar = \frac{4}{30\,000} \approx 1{,}33 \times 10^{-4}
    \label{eq:cell_ratio}
\end{equation}
ist prinzipiell messbar: In Präzisionsexperimenten an der Grenze der
Planck-Auflösung sollte die diskrete Zeitstruktur bei $T_0$, nicht bei $t_P$
sichtbar werden.

\subsection{Vorhersage 2: Automatenzustände und Teilchenmassen}

Wenn der Zellularautomat ein Grenzfall von T0 ist, dann müssen die
Energie-Eigenwerte des Automaten exakt dem T0-Massenspektrum entsprechen:
\begin{equation}
    E_n^{\text{Automat}} = E_0 \cdot \xipar^{-n}, \quad n = 0, 1, 2, \ldots
    \label{eq:automaton_mass}
\end{equation}
Dies ist eine direkte Verbindung zwischen der CAI-Struktur und experimentell
gemessenen Teilchenmassen.

\subsection{Vorhersage 3: Bellsche Korrelationen und $\xipar$-Korrekturen}

In der CAI ist Superdeterminismus quantitativ nicht ausgearbeitet. T0 liefert
eine konkrete Vorhersage: Bellsche Korrelationen sollten auf der Skala
$\xipar \approx 10^{-4}$ von den idealen Quantenvorhersagen abweichen:
\begin{equation}
    \langle AB \rangle_{T0} = \langle AB \rangle_{\text{QM}}
    \cdot \left(1 - \xipar \cdot \frac{\Delta x}{\ell_P}\right)
    \label{eq:bell_correction}
\end{equation}
wobei $\Delta x$ der räumliche Abstand der Messorte ist. Diese Korrektur ist
mit heutiger Technologie knapp unterhalb der Nachweisgrenze, aber prinzipiell
testbar.

\subsection{Vorhersage 4: Hamiltonoperator aus T0}

In der CAI bleibt der Hamiltonoperator ein Werkzeug ohne tiefen Ursprung. T0
sagt voraus, dass $H_{\text{op}}$ der zweiten Variation des Energiefunktionals
entspricht:
\begin{equation}
    H_{\text{op}} = \frac{\delta^2 E_{T0}}{\delta T^2}
    \bigg|_{T = T_{\text{vac}}}
    \label{eq:hamilton_from_t0}
\end{equation}
Dies ist eine testbare Aussage: Die Spektren von Quantensystemen sollten aus
dem T0-Zeitfeld-Energiefunktional berechnet werden können — ohne freie Parameter.

\clearpage

% ==============================================================================

\begin{figure}[htbp]
  \centering
  \includegraphics[width=\textwidth]{../fig/fig6_agreements_de}
  \caption{\"{U}bereinstimmungen (gr\"{u}n, links) und fundamentale Unterschiede (blau, rechts) zwischen CAI und T0-Theorie im \"{U}berblick.}
  \label{fig:agreements_de}
\end{figure}

\section{Gesamtvergleich: CAI und T0}
\label{sec:full_comparison}
% ==============================================================================

\begin{table}[h]
    \centering
    \caption{Systematischer Vergleich: CAI ('t Hooft 2016) und T0-Theorie}
    \label{tab:full_comparison}
    \resizebox{\textwidth}{!}{\begin{tabular}{@{}lll@{}}
        \toprule
        Aspekt & CAI ('t Hooft) & T0-Theorie \\
        \midrule
        \textbf{Grundaussage} & QM ist effektive Theorie eines CA & QM emergiert aus Zeitfeld-Geometrie \\
        \textbf{Determinismus} & Ja (ontologische Zustände) & Ja ($T(x,t)$ deterministisch) \\
        \textbf{Probabilismus} & Epistemisch & Epistemisch \\
        \textbf{Wellenpaketkollaps} & Kein echter Kollaps & Kein echter Kollaps \\
        \textbf{Hilbert-Raum} & Werkzeug & Werkzeug (Projektion des Zeitfeldes) \\
        \textbf{Zeitdiskretheit} & Postuliert ($\delta t \sim t_P$) & Abgeleitet ($T_0 = \xipar \cdot t_P$) \\
        \textbf{Zellenstruktur} & Postuliert & Geometrisch determiniert \\
        \textbf{Bell-Theorem} & Superdeterminismus (qualitativ) & Geometrische Korrelationen ($\xipar$) \\
        \textbf{Freie Parameter} & Nicht erklärt & Null (alle aus $\xipar$) \\
        \textbf{Fermionen} & Unvollständig & Topologische Windungszahlen \\
        \textbf{Gravitation} & Offen & Aus Zeitfeld-Gradient \\
        \textbf{Zeitpfeil} & Informationsverlust & Geometrisch erzwungen \\
        \textbf{Schwarze Löcher} & Informationsparadoxon offen & $T_{\min} = \xipar \cdot t_P \neq 0$, unitär \\
        \textbf{Standardmodell} & 20+ freie Parameter offen & Alle aus $\xipar$ \\
        \textbf{Status} & Programm, noch unvollständig & Vollständiger geometrischer Rahmen \\
        \bottomrule
    \end{tabular}}
\end{table}

\clearpage

% ==============================================================================
\section{Zusammenfassung}
\label{sec:summary}
% ==============================================================================

Gerard 't Hooft und die T0-Theorie teilen eine fundamentale Überzeugung: Die
Quantenmechanik ist keine ultimative Theorie, sondern die effektive Beschreibung
eines tieferliegenden deterministischen Mechanismus. Beide lehnen den
Wellenpaketkollaps, den fundamentalen Probabilismus und die ontologische
Rolle des Hilbert-Raumes ab.

In sechs fundamentalen Punkten stimmen CAI und T0 überein: Determinismus,
Ablehnung des Kollapses, Hilbert-Raum als Werkzeug, diskrete Zeitstruktur
auf der Planck-Skala, Quantenmechanik als Grenzfall klassischer Dynamik,
und Zeitpfeil als emergentes Phänomen.

In vier Punkten geht T0 fundamentaler: Statt den Automaten zu postulieren,
leitet T0 die diskrete Raumzeitstruktur aus $\xipar = 4/30\,000$ ab. Die
freien Parameter des Standardmodells emergieren aus $\xipar$. Gravitation
folgt aus dem Zeitfeldgradienten. Fermionen entstehen aus topologischen
Windungszahlen.

Die zentrale These: Der 't Hooftsche Zellularautomat ist ein Grenzfall der
T0-Geometrie — die diskrete Gitterapproximation des Zeitfeldes $T(x,t)$ auf
der Sub-Planck-Skala $L_0 = \xipar \cdot \ell_P$. 't Hoofts Programm zeigt,
dass deterministische Quantenmechanik \emph{möglich} ist. T0 zeigt, warum
sie \emph{notwendig} ist und wie sie konkret aussieht.

\begin{philosophical}[Einsteins letztes Wort]
    Einstein hatte recht: Gott würfelt nicht. Aber er spielte nicht Schach —
    er löste eine geometrische Variationsaufgabe. Das Minimum von
    $E_{T0} = \int |\nabla T|^2 d^4x$ unter der Nebenbedingung $T \cdot m = 1$
    ist das Universum.
\end{philosophical}

\begin{thebibliography}{99}

\bibitem{thooft_2016}
G.\ 't Hooft,
\newblock \emph{The Cellular Automaton Interpretation of Quantum Mechanics},
\newblock Fundamental Theories of Physics 185,
\newblock Springer, Cham, 2016.
\newblock DOI: 10.1007/978-3-319-41285-6.

\bibitem{thooft_arxiv}
G.\ 't Hooft,
\newblock The Cellular Automaton Interpretation of Quantum Mechanics,
\newblock arXiv:1405.1548 [quant-ph] (2014, rev.\ 2015).

\bibitem{pascher_xi}
J.\ Pascher,
\newblock Der geometrische Ursprungsparameter $\xi = 4/30\,000$,
\newblock T0-Dokumentationsreihe, Dokument 009 (2025).

\bibitem{pascher_grundlagen}
J.\ Pascher,
\newblock T0-Theorie: Zeit-Masse-Dualität als fundamentales Prinzip,
\newblock T0-Dokumentationsreihe, Dokument 003 (2025).

\bibitem{pascher_zeitpfeil}
J.\ Pascher,
\newblock Zeitpfeil als geometrische Konsequenz der Raumzeitstruktur,
\newblock T0-Dokumentationsreihe, Dokument 157 (2026).

\bibitem{pascher_qm_opt}
J.\ Pascher,
\newblock T0-Quantenmechanik: Geometrischer Formalismus und Optimierungsstrategien,
\newblock T0-Dokumentationsreihe, Dokument 034 (2025).

\bibitem{pascher_ising}
J.\ Pascher,
\newblock T0-Theorie und Coherent Ising Machines,
\newblock T0-Dokumentationsreihe, Dokument 161 (2026).

\bibitem{pascher_uv}
J.\ Pascher,
\newblock Ultraviolett-Katastrophe, Plancks Quantisierung und Einsteins Lichtquanten,
\newblock T0-Dokumentationsreihe, Dokument 163 (2026).

\bibitem{bell_1964}
J.\ S.\ Bell,
\newblock On the Einstein Podolsky Rosen paradox,
\newblock \emph{Physics Physique Fizika} \textbf{1}, 195--200 (1964).

\bibitem{bohm_1952}
D.\ Bohm,
\newblock A suggested interpretation of the quantum theory in terms of
hidden variables I \& II,
\newblock \emph{Physical Review} \textbf{85}, 166--193 (1952).

\end{thebibliography}
